% Standalone Overleaf document: compile with pdfLaTeX; no external assets.
% Preliminary plan agreed while Stage 3 was still running. Not an execution record.
\documentclass[11pt,a4paper]{article}
\usepackage[margin=2.2cm]{geometry}
\usepackage[T1]{fontenc}
\usepackage[utf8]{inputenc}
\usepackage{amsmath,amssymb,booktabs,enumitem}
\usepackage[hidelinks]{hyperref}
\setlength{\parindent}{0pt}
\setlength{\parskip}{5pt}
\setlength{\emergencystretch}{2em}
\setlist{itemsep=2pt,topsep=3pt}
\newcommand{\file}[1]{\nolinkurl{#1}}
\newcommand{\purpose}{\textbf{Question and rationale.} }
\newcommand{\procedure}{\textbf{Procedure.} }
\newcommand{\interpretation}{\textbf{Interpretation and output.} }
\title{Stage 4: Circuits and Inter-Head Communication\\\large Preliminary experimental plan}
\author{Phase A: Single-hop SIH Audit}
\date{22 September 2026}

\begin{document}
\maketitle

\section{Objective and scope}
Stage 3 characterizes individual heads. Stage 4 asks how those heads exchange
information, cooperate or compensate for one another, and form mechanisms that
support the model's single-hop \emph{mother-of} behavior. It also asks whether
the heads selected by the implemented Relation Index (RI) participate in these
mechanisms, including through effects hidden in single-head interventions.

This is a planning document, not a report of completed experiments. Stage-3
results are pending. Writer, reader and mover labels remain hypotheses until
supported by the corresponding measurements. This plan replaces the circuit
and group-analysis portions of the original methodology's Stages 3--4; the
revised Stage-3 plan remains the authority for individual-head characterization.

All discovery uses the existing 89 discovery families. The 87 held-out families
are new examples of the \emph{same mother-of task}, not a new semantic relation.
They are used only after the circuit and evaluation protocol are frozen.
The location world, a new observational metric and parent-coverage analysis are
excluded from this stage. We do not assume that one circuit explains all cases,
or that every interesting head belongs to an identifiable circuit.

\section{Shared definitions and decisions to freeze}
\subsection{Outcome, units and populations}
Retain the pinned OLMo-2-1124-7B checkpoint, tokenizer and Stage-2 metric.
For a clean/corrupted pair with first divergent answer tokens $y_c,y_r$ and
shared answer prefix $c$, define
\[
 M(x)=z(y_c\mid x,c)-z(y_r\mid x,c),\qquad
 I_T=M(x_c)-M(x_c;T).
\]
Here $T$ denotes the specified corrupted-donor or ablation intervention;
positive $I_T$ means movement away from the clean answer. Reverse recovery is
$J_T=M(x_r;T_{c\to r})-M(x_r)$, with the same answer sign.
Save both answer logits, not only their difference. Neither quantity is an
accuracy change or an additive share of the model's computation.

Use the original common 20 discovery families, both orders (40 pairs), for
screen calibration and initial exact tests. Extend retained edge and group
tests to all 89 discovery families, both orders, before freezing the circuit.
New query-control pairs are a separately reported contrast. Average orders
within family, then families equally. Report signed means, absolute-effect
summaries, family variation and support. Proposed uncertainty reporting is
20,000 whole-family bootstrap draws with a recorded seed; discovery intervals
are descriptive and hold the selected components fixed.

\subsection{Search space and circuit membership}
Start from the 105 Stage-3 heads, including the unchanged 25 random references.
Let $R$ denote the 59 test-only RI candidates plus L3H11/L9H22; these are
\emph{RI-selected candidates}, not established semantic heads. Keep the
historical pair and the test-only selection distinguishable in reporting.

The inventory is a starting point, not a closed universe. Screen incoming and
outgoing connections of seed heads against all temporally valid heads and MLP
blocks, including input and output connections. Expand around newly retained
components in documented rounds. A bounded search must report its unexplored
boundary; absence from it is not evidence of irrelevance.

A circuit node is a component at a declared token position or annotated span;
an edge additionally identifies the receiver channel (Q, K, V, MLP input or
output readout). Record exact token positions and their cross-example roles.
Candidate sites from Stage 3 guide priority, but do not exclude sites solely
because their single-head effect was small.

\textbf{Provisional membership rule.} A supported member must have (i) an
exactly tested causal route toward the answer, possibly conditional on a
specified group intervention, and (ii) a measurable conditional contribution
when removed from or restored to the proposed mechanism. A large group effect
does not establish every member's contribution. Preserve negative as well as
positive contributors. Distinguish supported members, candidates, unresolved
cases, and no detected contribution within the tested scope. A head can belong
to several mechanisms; passing these tests does not establish semantic
specificity by itself.

\subsection{Interventions and implementation contract}
\textbf{Two declared baselines.} Use paired corrupted-donor patching for the
primary edge/path and recovery experiments. For circuit retention/removal,
use discovery-mean ablation as the primary baseline, preserving the original
methodology, and paired-donor replacement as a sensitivity analysis. Estimate
means over discovery families only, balancing orders and base/corrupted
variants, matched by component and position role/token offset. Specify handling
of variable-length spans before use; never condition means on the correct
answer identity. Record the exact activation site where each mean is applied
and keep these means fixed for validation \cite{patchguide,bestpractices}.

\textbf{OLMo-specific requirement.} Attention outputs undergo shared
post-attention normalization. In general,
$N(S+\delta)-N(S)\ne N(\delta)$, where $S$ is the sum of projected head
outputs. Do not invent additive normalized head contributions. The graph must
retain shared normalization operations, actual pre-projection normalization,
QK normalization and RoPE. An adapted EAP score is the derivative of the
\emph{implemented intervention}, checked against finite differences and exact
replacement, not a weight-only formula assumed to transfer from GPT-2.

\textbf{Scope and gates.} Distinguish the final original colon from the metric
site after a shared answer prefix. Initial targeted interventions use original
prompt positions, with answer-prefix positions live, matching Stage 2. A full
circuit claim must additionally specify retained/ablated components at the
metric-prefix positions; otherwise label it a prompt-position subcircuit with
live background computation. Likewise, keeping all MLPs live gives a
head subcircuit conditional on those MLPs, not a complete model circuit.
Before measurement require unchanged-hook and self-patch checks, baseline
replication, Q/K/V reconstruction checks, and full-retention equivalence. Freeze
numeric tolerances from the established runtime and a separate implementation
pilot; do not relax them after failures.

\textbf{Execution manifest, still to be finalized after Stage 3.} Freeze seed
sites/groups, edge-ranking and expansion budgets, calibration sample strata,
screen-acceptance criteria, practically meaningful edge/member effect
thresholds, family-support rules, group subsets, ablation means, seeds and
stopping rules before their Stage-4 measurements. Do not transplant Stage-3
event-support thresholds or the head-level correlation gate to edges without
justification. Optional analyses need their own fixed site/layer/dimension
budgets. These unresolved fields make this a preliminary plan, not yet a
submission-ready protocol. Log discovery amendments before new batches;
validation cannot be used to revise them.

\section{Required discovery analyses}
The order follows dependencies. Early group tests run alongside edge mapping;
they need not wait for a finished graph. Reuse valid captures, but audit cache
contents and measure pilot runtime before describing work as CPU-only or cheap.

\subsection{Residual divergence across depth}
\purpose Locate where clean and corrupted representations differ, providing
a layer-level guide for targeted communication tests.

\procedure Capture the residual after each complete transformer block and at
the embedding boundary on the common 40 pairs. At the final original colon,
compute $d_\ell=\|r^c_\ell-r^r_\ell\|_2$. Report the corresponding residual
norms and relative distance
$d_\ell/[(\|r^c_\ell\|_2+\|r^r_\ell\|_2)/2]$; flag zero denominators.
Keep the post-prefix metric site separate where applicable. Plot family means
and variation. Once a candidate writer group exists, repeat the profile after
its intervention, with matched random-group controls.

\interpretation A jump between layers 11--16 is a hypothesis, not a constraint.
A flat distance can conceal rotations or transformations that make information
usable downstream. This curve cannot exclude important later heads or establish
a circuit edge. Deliver layer profiles and intervention-induced changes, with
exact capture boundaries and cache provenance.

\subsection{Position-aware edge screening and calibration}
\purpose Find plausible communication routes without exhaustively running
exact interventions on all edges. Position-aware scores reduce cancellation
across token roles \cite{eap,peap}.

\procedure For each allowed connection, estimate the effect of replacing its
source signal only along the declared receiver route. In an additive graph the
EAP approximation is $(z_u^r-z_u^c)^\top\nabla_{a_v}M$; use the corresponding
derivative through the actual OLMo intervention instead. Separate same-position
residual connections from cross-position attention computations: movement from
a fact token to the colon must traverse the receiver's attention operation.
Use semantic spans and token offsets to aggregate across examples and reorderings.

The usual two-forward/one-backward EAP cost is a conceptual baseline, not a
promise of identical runtime for this adaptation. Stage-2 saved head scores and
its $0.945$ head-ranking correlation do not validate edge estimates.
Calibrate exact effects on a preselected stratified sample spanning strong
positive/negative, middle, near-zero and random edge scores, channels, layers
and positions. Report sign agreement, magnitude error, ranking agreement and
strong effects missed by the screen. If inadequate under the frozen criteria,
test an explicitly defined EAP-IG variant or broaden exact sampling; neither
inherits the Stage-2 approximation's validation \cite{faith}.

\interpretation Deliver a ranked edge table with calibration evidence and
coverage. Retain both signs and random low-score checks; a low approximate score
is not grounds for a claim of no connection.

\subsection{Exact path patching and route expansion}
\purpose Determine whether a candidate source affects the answer through a
particular receiver channel, rather than merely having a correlated output
\cite{ioi,patchguide}.

\procedure For a candidate route from source $A$ to receiver $B$:
\begin{enumerate}
\item Cache clean and corrupted activations. In a clean recipient, replace
$A$ at the selected positions with its corrupted output. For a direct residual
route, freeze intervening attention-head and MLP branch outputs to clean
values, while keeping the designated normalization operations live. Other
routes require an explicitly listed set of live mediators.
\item Capture the resulting selected Q, K or V channel of $B$. Recompute
normalization and RoPE consistently; do not splice a pre-normalization vector
into a post-normalization site. Keep unselected channels and positions clean.
\item In a fresh clean run, insert only this hybrid receiver channel and let
the downstream model run normally. Compare its answer metric with the clean
baseline. The source itself is not patched in this final run.
\end{enumerate}
Save source/receiver positions, frozen and live components, replacement tensors
and answer logits. Test self-donors, matched wrong-position/channel controls and
reverse-direction recovery. Run retained routes on all 178 discovery pairs;
retain family/order heterogeneity. Test MLP-mediated routes separately where
screening suggests them, and expand around newly supported components under
the recorded budget.

\interpretation This measures the specified path intervention, not every
effect of $A$ on $B$. A failure on a direct route does not exclude an indirect
route. Shared normalization is part of the measured route, not an additive
contribution attributed solely to $A$. The proposed early-writer/late-reader
organization is one candidate outcome, not the search boundary.

\subsection{Fact-swap and query-swap contrasts}
\purpose Separate sensitivity to changed fact assignments from sensitivity
to which entity the question requests, without equating either with semantics.

\procedure On the common families, construct four conditions: original/swapped
mother assignments crossed with original/alternative query child. Reuse the
Stage-3 query controls where valid; generate the missing crossed condition.
Keep demonstrations fixed and check the gold answer and behavioral margin in
every condition. Do not filter for the desired attention pattern.
For each axis, compare matched endpoints differing only along that axis, using
their own original-versus-counterfactual answer sign. Recompute shared answer
prefixes. Patch only aligned sites: use verified span/anchor mappings for
query-length changes, and report unmatched sites rather than forcing alignment.
Apply the retained route and group tests to both axes and both orders; extend
retained contrasts to the full discovery set under a frozen selection rule.

\interpretation Report baseline task success, signed effects and recovery
relative to each contrast's gap, flagging small denominators. Importance on both
axes does not prove relational computation; importance on only one does not
prove lexical processing. Use reordering and consistent entity renaming as
targeted follow-ups when needed to separate location from identity. This is a
proposed diagnostic design, not a claim of methodological novelty.

\subsection{Group effects, backup roles and conditional membership}
\purpose Find cooperation, cancellation and backup contributions missed by
single-head averages, and test the participation of RI-selected candidates.

\procedure Define groups from Stage-3 profiles and verified routes, including
candidate writers/readers, opposite-sign contributors, and small-effect RI
candidates with plausible related roles. Preserve the actual grouping criteria.
For each group $G$, measure joint noising, mean ablation, and clean-donor
recovery in corrupted recipients at the declared sites. Plot effects against
coalition size using several seeded subsets/orders, with matched-size,
layer/site-matched random groups. Do not interpret one greedy curve as a
measure of distributedness.

For selected disjoint groups $A,B$, run baseline, $A$, $B$, and $A\cup B$ and
report $I_{A\cup B}-I_A-I_B$ under the same intervention and baseline. For a
candidate $h$, evaluate $I_{K\cup\{h\}}-I_K$ for declared backgrounds $K$
(including removal of plausible substitutes), and restoration of $h$ to the
corresponding reduced circuit. These tests can support a conditional role even
when $I_h$ alone is small. Restrict combinations by a recorded budget rather
than exhaustively testing all subsets.

For candidate writers, jointly restore clean activations at fact positions
only; compare the recovered answer with restoration at matched irrelevant
positions. This tests whether those activations support recovery given the
remaining corrupted model, not whether the writers alone solve the task.

\interpretation Deliver group curves, interaction tables and a conditional
contribution record for each proposed member. Non-additivity alone does not
identify its mechanism. No detected effect in sampled coalitions cannot rule
out every possible contribution. Analyze $R$ throughout this process, including
its small-effect members; do not postpone its assessment to a final overlap count.

\subsection{Circuit construction and discovery evaluation}
\purpose Assemble an explanation that reproduces the model's behavior and
test which components are required within the searched mechanisms.

\procedure Build input-to-output candidate graphs from verified routes, group
effects and required computational operations. Check discovery performance as
components are added or removed. Evaluate an edge-defined circuit by masking
excluded \emph{edges}, not just by ablating other heads while leaving all
connections among retained heads live. State the background treatment of
MLPs, normalization, embeddings, residual paths and answer-prefix sites.

For model or circuit $D$, define its same-sign clean--corrupted gap
\[
 G(D)=\frac1{89}\sum_f\frac12\sum_{o=0}^{1}
 \bigl[M_D(x^c_{f,o})-M_D(x^r_{f,o})\bigr],\qquad
 F(C)=\frac{G(C)}{G(\mathrm{full})}.
\]
The denominator uses the same families and orders and must be nonzero.
For donor-baseline sensitivity, replace excluded signals with the paired other
condition symmetrically and label this separate intervention. Use the primary
mean baseline for the inherited criterion $F(C)\geq0.80$. Also report
$|G(C)-G(\mathrm{full})|$, both condition means, per-family variation and answer
accuracy: an overshoot can reflect removal of negative contributors and is not
automatically a better explanation \cite{faith}.

\textbf{Completeness and minimality.} Remove the circuit from the full model,
but also compare $C\setminus K$ with $\mathrm{full}\setminus K$ for predefined
role groups and seeded subsets $K\subseteq C$. A large discrepancy reveals
possible missing backup computation. For each claimed member, seek a declared
background in which its removal/restoration changes performance; apply the same
logic to claimed edges. Report tested subsets and search limits. A greedy
reduction establishes minimality only relative to that procedure, not a unique
globally minimal circuit \cite{ioi}.

\interpretation Deliver circuit-size/performance curves, the retained graph,
ablation sensitivity and any unresolved completeness failures. Return to
discovery expansion when justified and budgeted; do not open validation to
choose a more favorable circuit.

\section{Conditional interpretation analyses}
These deepen an established causal finding; neither is required for every
head or edge. Run them on discovery families only, with bounded sweeps and
explicitly recorded triggers.

\subsection{Communication-channel content}
\purpose Explain what an important verified route carries: candidate identity,
position-related information or another signal.

\procedure Trigger this analysis when a verified route's content remains
unexplained. Decompose the writer's $W_{OV}=\sum_i\sigma_i u_i v_i^\top$,
not a writer--reader product. Use its components as candidate directions;
weight-composition scores are descriptive screening only. Compare component
removal/restoration along the verified route with controls matched for rank
and actual intervention magnitude, sampling random directions in the same
writer output space. Preserve real value inputs, attention and shared
normalization, and include complementary-subspace controls. A global weight
edit affects other routes and cannot by itself establish this channel
\cite{talking,framework}.

Use minimal contrasts that change identity at fixed sites and move the same
entities across sites, with tokenization and answer annotations checked.
Evaluate whether the channel's effect follows identity or position and whether
it changes the receiver's routing or retrieved content. Neither outcome must
fit a binary pointer/content label.

\interpretation Report a supported subspace only if its route-specific causal
effects exceed the matched controls under the frozen criterion. Singular value
size and vocabulary projections alone do not establish what information the
model uses.

\subsection{Controlled representation readouts: Patchscopes-style}
\purpose Read information at a causally localized site when direct vocabulary
projections do not explain it, particularly before/after a group intervention
\cite{patchscopes}. Reuse an adequate Stage-3 optional readout if already done.

\procedure For writer sites, inject the residual into a fixed token-identity
few-shot readout. For receiver inputs, also test a fixed relation-completion
readout. Bound source/target-layer sweeps in advance and preserve residual
boundary conventions. Use generic demonstration names disjoint from the test
names; any supplied query entity must be identical across matched readouts.
Compare clean, corrupted and head/group-intervened representations in the same
readout prompt, alongside no injection, irrelevant injection and entity-swap
controls. Save all sweep results, not only readable examples.

For each prespecified candidate name, report teacher-forced sequence log
probability and candidate ranking. For generated text, the SelfIE-style
relevance diagnostic compares each token's probability with and without
injection, conditioned on the \emph{same generated prefix} \cite{selfie}.
Treat free text as spot-check evidence; check whether intervention changes the
decoded candidate consistently across families.

\interpretation Residuals contain many components, and the decoder prompt can
perform additional computation. Fictitious names do not exclude lexical copying
or prompt cues. Injection relevance measures an effect on the readout, not the
truth of its explanation. Successful decoding does not establish original-task
use; failed decoding does not establish information absence. Attribution to a
group requires the matched group-intervention comparison and supporting causal
tests.

\section{Frozen validation and final RI assessment}
\purpose Test whether the selected explanation transfers to unseen families
of the same task, and answer where RI-selected heads participate.

\procedure Before opening the 87 families, archive the circuit(s), membership
and edge rules, masks, mean activations, metric definitions, selected group/path
tests, optional-analysis decisions, code/input hashes and statistical plan.
Evaluate baseline behavior and the frozen circuit under the primary ablation
and prespecified sensitivity condition. Apply the gap formula above with 87
families, retaining both orders; report the $80\%$ criterion, uncertainty and
condition-specific behavior. Do not discard validation families because the
circuit fails on them. ``One evaluation'' means one frozen evaluation suite,
not one forward pass. Any revision after seeing its results is exploratory and
requires fresh data for another confirmatory claim.

For $R$, report both $|R\cap C_H|/|R|$ and $|R\cap C_H|/|C_H|$, where $C_H$
is the set of supported head members; state if the search incompletely covers
$R$. Preserve its historical/test-only split. Add a head-level table of verified
routes, conditional group contributions, sign, token roles and validation
status. Compare descriptive overlap with layer-matched reference sets, but
base mechanistic conclusions on interventions. A candidate absent from one
compact faithful circuit may still act in another valid or backup mechanism.

\section{Deliverables and completion}
Provide (i) a frozen execution manifest, (ii) residual and calibrated edge
tables, (iii) exact route/group and conditional-membership evidence,
(iv) candidate circuit graphs with discovery/validation faithfulness and
completeness diagnostics, and (v) the RI participation table. Each test is
marked completed, insufficient support, unavailable input or unresolved
implementation; optional tests are additionally marked not triggered when
appropriate. Missing measurements are never zero-filled.

Stage 4 is complete when the declared search and evaluation scope is executed
and reported, not when every head receives a label or the $80\%$ criterion is
met. A failed criterion, several overlapping mechanisms or a partially explained
circuit are legitimate outcomes. No claim extends automatically to other
relations, two-hop reasoning or all circuits in the model.

\begin{thebibliography}{99}
\bibitem{eap} A. Syed, C. Rager and A. Conmy (2023).
\emph{Attribution Patching Outperforms Automated Circuit Discovery}.
\url{https://arxiv.org/abs/2310.10348}
\bibitem{peap} T. Haklay et al. (2025).
\emph{Position-aware Automatic Circuit Discovery}.
\url{https://arxiv.org/abs/2502.04577}
\bibitem{ioi} K. Wang et al. (2022).
\emph{Interpretability in the Wild: a Circuit for Indirect Object Identification in GPT-2 Small}.
\url{https://arxiv.org/abs/2211.00593}
\bibitem{patchguide} S. Heimersheim and N. Nanda (2024).
\emph{How to Use and Interpret Activation Patching}.
\url{https://arxiv.org/abs/2404.15255}
\bibitem{bestpractices} F. Zhang and N. Nanda (2024).
\emph{Towards Best Practices of Activation Patching in Language Models: Metrics and Methods}.
\url{https://arxiv.org/abs/2309.16042}
\bibitem{faith} M. Hanna, S. Pezzelle and Y. Belinkov (2024).
\emph{Have Faith in Faithfulness: Going Beyond Circuit Overlap When Finding Model Mechanisms}.
\url{https://arxiv.org/abs/2403.17806}
\bibitem{talking} J. Merullo, C. Eickhoff and E. Pavlick (2024).
\emph{Talking Heads: Understanding Inter-layer Communication in Transformer Language Models}.
\url{https://arxiv.org/abs/2406.09519}
\bibitem{framework} N. Elhage et al. (2021).
\emph{A Mathematical Framework for Transformer Circuits}.
\url{https://transformer-circuits.pub/2021/framework/index.html}
\bibitem{patchscopes} A. Ghandeharioun et al. (2024).
\emph{Patchscopes: A Unifying Framework for Inspecting Hidden Representations of Language Models}.
\url{https://arxiv.org/abs/2401.06102}
\bibitem{selfie} H. Chen, C. Vondrick and C. Mao (2024).
\emph{SelfIE: Self-Interpretation of Large Language Model Embeddings}.
\url{https://arxiv.org/abs/2403.10949}
\end{thebibliography}
\end{document}
